\documentclass[11pt,a4paper]{article}
\usepackage[utf8]{inputenc}
\usepackage[T1]{fontenc}
\usepackage{lmodern}
\usepackage{amsmath,amssymb,amsfonts,amsthm}
\usepackage{geometry}
\geometry{margin=1in}
\usepackage{xcolor}
\usepackage{hyperref}
\usepackage{cleveref}
\usepackage{booktabs}
\usepackage{array}
\usepackage{microtype}
\usepackage{tcolorbox}

\hypersetup{
    colorlinks=true,
    linkcolor=blue,
    citecolor=red,
    urlcolor=cyan,
    pdftitle={Certified Lattice Data and a Within-Family Selector for the Cooper s7/s10 K3 Families},
    pdfauthor={Xavier Callens and the SocrateAI Scientific Agora Collaboration}
}

\newtheorem{theorem}{Theorem}[section]
\newtheorem{proposition}[theorem]{Proposition}
\newtheorem{definition}[theorem]{Definition}
\theoremstyle{remark}
\newtheorem{remark}[theorem]{Remark}

% Epistemic tier badges, matching Paper 12's convention.
\newcommand{\tierA}{\colorbox{green!15}{\textsf{\scriptsize Tier A --- Lean kernel}}}
\newcommand{\tierB}{\colorbox{yellow!20}{\textsf{\scriptsize Tier B --- derived, not measured}}}
\newcommand{\tierL}{\colorbox{blue!10}{\textsf{\scriptsize Tier L --- literature}}}
\newcommand{\tierC}{\colorbox{orange!15}{\textsf{\scriptsize Tier C --- conjecture}}}
\newcommand{\lean}[1]{\texttt{#1}}

\newtcolorbox{headlinebox}{%
  colback=red!4, colframe=red!45!black, boxrule=0.5pt, arc=1pt,
  left=6pt, right=6pt, top=4pt, bottom=4pt,
  fonttitle=\bfseries\small, title={Headline result}
}
\newtcolorbox{scopebox}{%
  colback=gray!5, colframe=gray!45, boxrule=0.4pt, arc=1pt,
  left=6pt, right=6pt, top=3pt, bottom=3pt,
  fonttitle=\bfseries\small, title={What is not claimed}
}
\newtcolorbox{correctionbox}{%
  colback=orange!6, colframe=orange!55!black, boxrule=0.5pt, arc=1pt,
  left=6pt, right=6pt, top=4pt, bottom=4pt,
  fonttitle=\bfseries\small, title={Corrections recorded in this session}
}

\title{\Large \textbf{Certified Lattice Data and a Within-Family Selector}\\[0.2em]
\textbf{for the Cooper $s_7$/$s_{10}$ $K3$ Families}}

\author{\textbf{Xavier Callens} \and \textbf{SocrateAI Scientific Agora Collaboration} \\[0.3em]
\textit{Stream 2 --- Selection \& Geometry (\texttt{SocrateAI-Scientific-Agora-K3-DarkMatter})}}
\date{2026-09-29 \quad (preprint; not reviewed)}

\begin{document}
\maketitle

\begin{scopebox}
This is a mathematics note about lattice arithmetic on two families of $K3$ surfaces. It makes no
physical claim of any kind: no dark-matter identification, no dark-energy statement, no observable.
The programme's flux/tadpole route is blocked pending a specified threefold base (F5b) and a
selection among the surfaces discussed here maps to no measurement. These renunciations are load-bearing,
not boilerplate; see \cref{sec:scope}.
\end{scopebox}

\begin{abstract}
We report a session's results on the Cooper $s_7$ and $s_{10}$ Picard--Fuchs families, all of it lattice
mathematics with an explicit epistemic tier on every claim. Six points of the $\rho=20$ locus of each
family's transcendental lattice are exhibited as exact vectors with saturated orthogonal complements,
reduced binary forms, and frame indices; the lattice arithmetic of all twelve is kernel-checked in Lean~4,
three of the six \emph{independently} by two separate developments, which we re-verify ourselves before
citing either. A within-family selector --- minimal discriminant of the transcendental lattice --- is
adopted by an explicit ruling among two computed alternatives, chosen because the rejected alternative is
not unique in one family. An explicit $M_7$-polarized model realizes the $s_7$ family as an Inose surface;
its elliptic fibration, taken from a sourced literature equation and re-derived here, has the two
singular loci of the Picard--Fuchs operator indistinguishable at the level of discriminant-root orders,
which answers a standing question about one of them in the negative. Stage-2 monodromy of the order-two
operator underlying each family is certified by rigorous ball arithmetic, replacing an earlier numerical-tolerance
heuristic with enclosures whose diameters are reported exactly. We record, in band, five errors made and
caught during the same session, because the standing discipline of this programme requires retractions to
be machine-visible rather than a footnote.
\end{abstract}

\tableofcontents

\section{Scope and what is not claimed}\label{sec:scope}

\begin{scopebox}
No claim in this note is physical. The $\rho=20$ locus is a set of lattice points on a modular curve; it
maps to no observable and licenses no dark-sector identification. The discriminant-floor selector of
\cref{sec:selector} is a naming convention for which candidate a certificate calls ``distinguished''
within a family, adopted by an explicit ruling, and it does not rank one Cooper family against the other.
No Kodaira fibre type is attached to any point of either family's Picard--Fuchs operator: that reading is
a category error for this class of family (elliptic points of a modular curve are not Kodaira
degenerations), and the one place an explicit Weierstrass model is computed (\cref{sec:fibration}) we
report discriminant-root \emph{orders} only, withholding the label pending a separate ruling. The
$s_{10}$ family's transcendental lattice remains a DRAFT certificate; every $s_{10}$ result here is
advisory and does not depend on, or promote, that certificate's status.
\end{scopebox}

\section{The six $\rho=20$ locus rows, kernel-checked}\label{sec:rows}

Each register family ($s_7$ at level $N=7$, $s_{10}$ at level $N=10$) carries a rank-three lattice
$U\oplus\langle 2N\rangle$ with quadratic form $v^2 = 2xy+2Nz^2$ on coordinates $v=(x,y,z)$. Three points
per family are singled out by the vanishing locus of the family's order-three Picard--Fuchs operator
$L_3$: the two finite singular points and the point at infinity. At each such point the transcendental
lattice, rank $19$ generically, drops to rank $2$ and becomes positive definite --- $\rho$ jumps to $20$.
\Cref{tab:rows} lists all six, as exact vectors with their norm, the exact orthogonal complement (its Gram
matrix, reduced to a binary quadratic form of discriminant $D$), and the Lean kernel status.

\begin{table}[h]
\centering
\begin{tabular}{lllrrll}
\toprule
Family & $z$ & $v$ & $v^2$ & $D$ & $T$ (reduced) & Tier / sources \\
\midrule
\lean{cooper\_s7} & $-1$ & $(2,-4,1)$ & $-2$ & $-7$ & $[2,1;1,4]$ & Tier A, 2 source(s) \\
\lean{cooper\_s7} & $1/27$ & $(1,-1,0)$ & $-2$ & $-28$ & $[2,0;0,14]$ & Tier A, 2 source(s) \\
\lean{cooper\_s7} & $\infty$ & $(14,-14,-5)$ & $-42$ & $-3$ & $[2,1;1,2]$ & Tier A, 2 source(s) \\
\lean{cooper\_s10}\,(ADVISORY) & $-1/4$ & $(2,-6,1)$ & $-4$ & $-20$ & $[4,2;2,6]$ & Tier A, 1 source(s) \\
\lean{cooper\_s10}\,(ADVISORY) & $1/16$ & $(1,-1,0)$ & $-2$ & $-40$ & $[2,0;0,20]$ & Tier A, 1 source(s) \\
\lean{cooper\_s10}\,(ADVISORY) & $\infty$ & $(10,-10,-3)$ & $-20$ & $-4$ & $[2,0;0,2]$ & Tier A, 2 source(s) \\
\bottomrule
\end{tabular}

\caption{The six locus rows. \tierA\ for the lattice arithmetic (norm, saturated complement, Gram, reduced
form, frame determinant) of every row; three rows carry two independent kernel developments. \tierB: the
numeric recognition of $z$, and the identification of the orthogonal complement with the family's
transcendental lattice $T_X$ (Dolgachev 1996 \S7, Doran 1998 Thm.~5.13, \tierL). The index step
$|\det(\text{frame})| = |v^2|/\operatorname{div}(v)$ is checked per row, never proved in general, in
either kernel development.}\label{tab:rows}
\end{table}

\begin{headlinebox}
For the two rows at $z=1/27$ and $z=-1$ (the $s_7$ family) and, independently added this session, both
families' $z=\infty$ rows, the lattice half is kernel-checked \emph{twice}, by two developments written
without reading each other's source, then read at the statement level and re-gated by us in each
producer's own build tree before either is cited.
\end{headlinebox}

\begin{remark}
The two independent developments state the norm, the exact orthogonal complement as a set (an
if-and-only-if, not merely two orthogonal vectors), the complement's Gram matrix, its Gauss-reduced form,
and the frame determinant, each as a standalone kernel-checked declaration, over the standard axiom set
$\{\texttt{propext},\ \texttt{Classical.choice},\ \texttt{Quot.sound}\}$ and nothing else. Producer and
verifier are different in every case reported here: we re-ran the build, the sorry grep, the axiom audit,
the statement lock, and \texttt{\#print axioms} ourselves, in the producer's own worktree, before writing
any lattice tier into a certificate.
\end{remark}

\section{An adopted selector, and the alternative it rejected}\label{sec:selector}

Six candidates are too many to call ``the'' surface of a family without saying which quantity is being
extremized. Two computed candidates were compared:

\begin{itemize}
\item \textbf{SEL-D}: minimal $|\operatorname{disc}T|$ within a family, verified unique not merely by
  class number (which counts reduced forms) but by the count of points \emph{on the modular curve}
  itself.
\item \textbf{SEL-N}: minimal $|v^2|$ within a family --- the $-2$ floor achieved by an $N$-uniform Lean
  theorem in both kernel developments.
\end{itemize}

\Cref{tab:selector} gives both, exactly, per family.

\begin{table}[h]
\centering
\begin{tabular}{llll}
\toprule
Family & SEL-D (min $|{\rm disc}\,T|$) & SEL-N (min $|v^2|$) & Agree? \\
\midrule
\lean{cooper\_s7} & $T=[2,1;1,2]$, $D=-3$ & tie: $D=-7$, $D=-28$ (2 rows) & No \\
\lean{cooper\_s10} & $T=[2,0;0,2]$, $D=-4$ & $D=-40$ (unique) & No \\
\bottomrule
\end{tabular}

\caption{SEL-D and SEL-N disagree in both families. SEL-N is \emph{not unique} for $s_7$: two rows tie at
the $-2$ floor, corresponding to the two fixed points of the Atkin--Lehner involution $W_7$ on the
family's Hauptmodul coordinate.}\label{tab:selector}
\end{table}

\textbf{SEL-D was adopted.} The choice was put, with both alternatives computed and the disagreement
stated, to the programme's designated decision authority, who selected SEL-D directly. This narrows an
earlier, separately-ruled clause (``no minimum-discriminant rule'') to cross-family comparisons only; it
does not rank $s_7$ against $s_{10}$, and it does not promote $s_{10}$'s DRAFT lattice certificate. The
selected surfaces --- $T=A_2$ ($D=-3$) for $s_7$, $T=\langle2\rangle\oplus\langle2\rangle$ ($D=-4$) for
$s_{10}$, advisory --- extremize the same quantity that a separate, pre-existing attractor-mechanism
development in this programme's Lean corpus already names for the smallest admissible black-hole charge
configuration; both are instances of the same discriminant-gap floor ($4ac-b^2\ge3$), so the coincidence
is forced by that shared floor and is stated once, here, as an observation and not as corroborating
evidence for either result.

\section{An explicit $M_7$-polarized model}\label{sec:model}

Restricting the Clingher--Doran--Lewis--Whitcher normal form for $M$-polarized $K3$ surfaces to the
$M_7$-polarized locus realizes the $s_7$ family as an Inose surface $X(a(z),b(z))$, with $a^3=\pi(z)$,
$b^2=\pi(z)-\sigma(z)+1$ and
\[
\sigma(z) = J(\tau)+J(7\tau), \qquad \pi(z) = J(\tau)\,J(7\tau),
\]
where $z$ is the family's own $\Gamma_0(7)^+$ Hauptmodul coordinate and $J$ is Klein's $J$-invariant,
normalized $J(i)=1$. Because $X_0(7)^+$ has a single cusp, $\sigma$ and $\pi$ are Laurent polynomials in
$z$; both were solved from their $q$-expansions and verified against every held-out coefficient to order
$q^{128}$ (\tierA\ for the finite-order verification; \tierB\ for the identification of $z$ with the
Hauptmodul coordinate). Exactly,
\[
\sigma(z)^2-4\pi(z) \;=\; \frac{-(z+1)(27z-1)(2z-1)^2(8z-1)^2(24z-1)^2(8z^2+16z-1)^2}{2985984\,z^{14}}.
\]
The \emph{simple} zeros of this expression are exactly the two finite singular loci of $L_3$
($z=-1$, $D=-7$; $z=1/27$, $D=-28$): the points where the two isogenous elliptic curves become
isomorphic, which is where $J(\tau)-J(7\tau)$ behaves like a square root, forced by the involution that
swaps the two curves. The double zeros are the other self-$7$-isogenous CM points in the same window,
each occurring exactly $h(D)$ times.

\subsection{The fibration}\label{sec:fibration}

A degree-two base change of a sourced elliptic pencil on $\mathrm{Km}(E_1\times E_2)$
(Kuwata--Shioda, \S5.3, \tierL) gives the Inose fibration explicitly, with the transcription checked
against the source's own printed discriminant before any further computation. \Cref{tab:fibration} gives
the discriminant-root orders --- never a Kodaira label, per \cref{sec:scope} --- for the generic case,
both $s_7$ singular loci, and the two special values $J=0,1$.

\begin{table}[h]
\centering
\begin{tabular}{lll}
\toprule
Case & Discriminant-root orders & Note \\
\midrule
Generic ($E_1\not\cong E_2$) & $\{10,10,1,1,1,1\}$ & --- \\
cooper\_s7, $z=-1$ & $\{10,10,2,1,1\}$ & $J=-125/64$ \\
cooper\_s7, $z=1/27$ & $\{10,10,2,1,1\}$ & $J=614125/64$ \\
$J=1$ ($E_1\cong E_2$) & $\{10,10,2,2\}$ & --- \\
$J=0$ ($E_i\cong E_\omega$) & $\{10,10,4\}$ & --- \\
\bottomrule
\end{tabular}

\caption{Discriminant-root orders of the Inose fibration. The mechanism behind the middle two rows: when
the two elliptic curves become isomorphic, the discriminant's constant term vanishes identically, and
after the base change the two simple roots collapse onto the branch point, producing exactly one
order-$2$ root there.}\label{tab:fibration}
\end{table}

\begin{headlinebox}
The two singular loci of $s_7$ are \emph{indistinguishable} at the level of discriminant-root orders:
both give $\{10,10,2,1,1\}$. A standing hypothesis in this programme --- that the two loci differ in
their fibre configuration --- is answered in the negative by this computation, for this fibration. The
second algebraic class responsible for the rank jump at either point must therefore be a Mordell--Weil
section, not a reducible fibre.
\end{headlinebox}

\section{Certified stage-two monodromy}\label{sec:monodromy}

The order-two Picard--Fuchs operator underlying each family's Sym$^2$ structure has monodromy that was
previously recognized only by a numeric-tolerance heuristic ($60$-digit precision, a $10^{-35}$ gate).
This session replaced that heuristic with rigorous ball arithmetic (Arb, via \texttt{python-flint}):
Frobenius series tails and Taylor-continuation tails are bounded by majorant induction on the operator's
own recurrence, giving enclosures whose diameters certify that each recognized monodromy entry is the
\emph{unique} rational of bounded denominator inside its enclosure. \Cref{tab:monodromy} reports the
diameters achieved; the certified matrices were then fed into the exact lattice pipeline and reproduce
every field of the corresponding lattice certificate, closing the chain from certified numerics to exact
lattice arithmetic.

\begin{table}[h]
\centering
\begin{tabular}{lll}
\toprule
Family & Clause & Enclosure diameter \\
\midrule
\lean{cooper\_s7} & $z=-1$ & $1.6\times 10^{-96}$ \\
\lean{cooper\_s7} & $z=1/27$ & $1.0\times 10^{-115}$ \\
\lean{cooper\_s10} & $z=-1/4$ & $4.2\times 10^{-102}$ \\
\lean{cooper\_s10} & $z=1/16$ & $2.2\times 10^{-115}$ \\
\bottomrule
\end{tabular}

\caption{Enclosure diameters of the certified stage-two monodromy. All are many orders of magnitude below
the denominator bound required for the recognized rational to be unique.}\label{tab:monodromy}
\end{table}

\section{Corrections recorded in band}\label{sec:corrections}

\begin{correctionbox}
Five errors were made and caught within this session, listed because the standing discipline of this
programme requires a retraction to be machine-visible, not a prose footnote.
\begin{enumerate}
\item An initial claim that SEL-N (\cref{sec:selector}) was unique and level-uniform was false: two rows
  of the $s_7$ family tie at the $-2$ norm floor. Caught before any recommendation was written down.
\item An initial reading of a standing ``no minimum-discriminant rule'' as scoped to cross-family
  comparison only was a stretch of the text, not what it said; the question was put to the decision
  authority rather than resolved by inference.
\item A transcription of the Inose (0,0)/(1,0) labels for two special points was swapped; corrected after
  an independent Weierstrass-reduction computation disagreed with it.
\item Two software defects were found while extending the Lean-attestation machinery of
  \cref{sec:rows}: a worktree-path resolver that only recorded one build tree per repository, so a
  second, later commit on a different branch of the same repository silently failed its own file-hash
  check; and a test control that matched a lattice row by vector coordinates alone, which collide across
  the two register families.
\item A version-control error landed a pull request containing only a brief document, omitting the
  actual code and data changes its own description claimed; caught by comparing a certificate's hash on
  the merged branch against the local file, and corrected by a follow-up commit before any downstream
  party was told the (wrong) hash.
\item \emph{Dated correction, 2026-10-07.} The sentences of \cref{sec:scope} and \cref{sec:selector} calling the
  $s_{10}$ lattice certificate a DRAFT, and the ``(ADVISORY)'' tags on the $s_{10}$ rows of
  \cref{tab:rows,tab:selector}, were true when written on 2026-09-29 and were superseded later that day by a
  ruling that made the value-identical certificate LIVE; the certificates were re-emitted without the flag on
  2026-10-07 and the two tables, which are generated from them, now show none. No value in any table changed.
  The prose sentences are left as written and are dated by this item; nothing in this paper ranks the two
  families, and that was never conditional on the flag.
\end{enumerate}
\end{correctionbox}

\section{Reproducibility}

Every number in \cref{tab:rows,tab:selector,tab:fibration,tab:monodromy} is generated directly from
certificates in \texttt{data/certificates/} by \texttt{scripts/render\_paper\_tables.py}, which fails
closed (\texttt{--check}) if any table fragment drifts from the certificates that produced it; nothing in
this document's tables is a typed literal. The checkers producing the underlying certificates each ship
negative controls that must fail on a tampered input before the corresponding positive result is trusted.
A full account of every certificate, checker, and the exact gate output for both independent Lean
developments of \cref{sec:rows} is in the repository's decision and status records, referenced in
\cref{sec:refs}.

\section*{Provenance}
Generated-by: Claude (Sonnet~5), Stream~2 \textbar\ Verified-by: \texttt{scripts/render\_paper\_tables.py
--check}; every cited checker's own controls \textbar\ Reviewed-by: N (preprint; not reviewed).

\section{References and source records}\label{sec:refs}

\begin{thebibliography}{9}
\bibitem{Dolgachev1996} I.~Dolgachev, ``Mirror symmetry for lattice polarized $K3$ surfaces,''
  \textit{arXiv:alg-geom/9502005} (1996).
\bibitem{Doran1998} C.~F.~Doran, ``Picard--Fuchs uniformization: modularity of the mirror map and
  mirror-moonshine,'' \textit{arXiv:math/9812162} (1998).
\bibitem{CD2007} A.~Clingher and C.~F.~Doran, ``Modular invariants for lattice polarized $K3$ surfaces,''
  \textit{arXiv:math/0602146} (2007).
\bibitem{CDLW} A.~Clingher, C.~F.~Doran, J.~Lewis, and U.~Whitcher, ``Normal forms, $K3$ surface moduli,
  and modular parametrizations,'' \textit{arXiv:0712.1880}.
\bibitem{DHNT} C.~F.~Doran, A.~Harder, A.~Y.~Novoseltsev, and A.~Thompson, ``Families of lattice polarized
  $K3$ surfaces with monodromy,'' \textit{arXiv:1312.6434}.
\bibitem{KS} M.~Kuwata and T.~Shioda, ``Elliptic parameters and defining equations for elliptic
  fibrations on a Kummer surface,'' \textit{arXiv:math/0609473}.
\bibitem{Moore1998} G.~Moore, ``Arithmetic and Attractors,'' \textit{arXiv:hep-th/9807087} (1998).
\bibitem{RepoRecords} SocrateAI Scientific Agora Collaboration, \texttt{SocrateAI-Scientific-Agora-K3-DarkMatter},
  \texttt{K3\_CRITERIA.md} and \texttt{TODO.md} at commit \texttt{d3a2a01} (2026-09-29); decision records
  \texttt{briefs/T0\_DECISIONS\_2026\_09\_28\_STREAM2.md} (selector ruling),
  \texttt{briefs/STREAM2\_AM8\_SELECTOR\_COMPARISON\_2026\_09\_28.md} (comparison detail),
  \texttt{briefs/WP\_GE10\_INOSE\_MODEL\_M7\_2026\_09\_27.md} (fibration derivation).
\end{thebibliography}

\end{document}
